\chapter[Logarithmic Linearization]{Logarithmic Linearization}\label{ch:sec:graphs}

The WIS factorisation $n_{mc}=\beta(c)/w(m)$ is multiplicative.
Geometry, however, is linear: distances, angles, and inner products all
require additive structure.  To bring geometric tools to bear on
encryption systems, we must linearise the WIS.  Taking logarithms
transforms the multiplicative factorisation into an additive linear
structure, revealing that the WIS representation is equivalent to a
linear condition on log-incidence vectors.

The logarithmic linearisation of this chapter is the bridge between algebra and geometry. It produces the mismatch operator $\Phi$, which will be the central object of Chapters~7 and~8, where it defines the least-squares problem and the posterior Laplacian.

which is the bridge from combinatorics to operator theory.  This
section shows that every WIS canonically determines a linear
representation whose associated operator is the central object of the
analytic theory.
\section{Logarithmic potentials}\label{ch:sec:logarithmic-potentials}
\begin{definition}[Logarithmic potentials]\label{ch:def:log-potentials}
For a WIS $\mathfrak W$, define vertex functions
$u(m)=\log w(m),\qquad v(c)=\log\beta(c).$
\end{definition}

\paragraph{Why this definition.}
The logarithmic potentials convert the multiplicative factorisation into an additive condition.  This linearisation opens the door to Hilbert-space geometry: additive structure means linear operators, least-squares, and spectral theory.

\begin{proposition}\label{ch:prop:potential}
For every incidence $mc$, $\log n_{mc}=v(c)-u(m).$ Consequently, the edge labelling $mc\mapsto\log n_{mc}$ is a difference of vertex potentials on the bipartite incidence graph.
\end{proposition}

\paragraph{Why logarithmic coordinates?}
The WIS factorisation $n_{mc}=\beta(c)/w(m)$ is multiplicative.  Geometry, however, is linear: distances, angles, and inner products all require additive structure.  Logarithms convert the multiplicative relation into the additive relation $\log n_{mc}=v(c)-u(m)$, making the problem amenable to linear algebra.  Why not use odds $P/(1-P)$ or the Fisher metric directly?  Odds are designed for binary hypothesis testing, not for the general incidence structure of encryption.  The Fisher metric is a Riemannian structure on the space of distributions, not a tool for factoring incidence matrices.  Logarithms are the unique transformation that turns the multiplicative ratio $\beta/w$ into a difference while preserving the order structure of probabilities.

\begin{proof}
Take logarithms in the defining identity $n_{mc}=\beta(c)/w(m)$.
\end{proof}

This identity is not merely a convenient rewriting.  It reveals that
the WIS data carry an additive structure that standard multiplicative
reasoning cannot access.  Because $u$ and $v$ are real-valued functions
on different vertex sets, they naturally belong to different vector
spaces, yet their difference lives on the common set of edges.  This
observation is the reason why Hilbert-space methods become available.
\section{Canonical liftings}\label{ch:sec:canonical-liftings}
Let $V_M=\mathbb R^{\Mset}$, $V_C=\mathbb R^{\Cset}$, and
$V_I=\mathbb R^I$, each with the standard Euclidean inner product.
Message potentials belong to $V_M$, ciphertext potentials to $V_C$;
both are lifted to the common incidence space $V_I$.
\begin{proposition}[Canonical lifting]\label{ch:prop:canonical-lifting}
The operators
$A:V_M\to V_I,\qquad (Au)_{(m,c)}=u_m,$ $C:V_C\to V_I,\qquad (Cv)_{(m,c)}=v_c,$ embed the potential spaces into $V_I$.  Both are injective when every
vertex has at least one incident edge.
\end{proposition}
\begin{proof}
If $Au=0$, then $u_m=0$ for every message that appears in an
incidence; if the incidence graph has no isolated vertices, this
covers all messages.  The same argument applies to $C$.
\end{proof}

\begin{quote}
\textbf{Mathematical intuition.}
The mismatch operator $\Phi=[A\; -C]$ sends a pair of vertex potentials $(u,v)$ to the vector $\Phi(u,v)$ whose components are $u_m-v_c$.  The logarithmic incidence vector is $\omega=-\Phi(u,v)$, equivalently $\omega=\Phi(-u,-v)$, with components $\omega_{mc}=v_c-u_m$.
Taking logarithms turns multiplicative relations into additive ones, enabling the use of linear algebra and Hilbert-space geometry.
\end{quote}

\section{Canonical linear representation}\label{ch:sec:canonical-linear-representation}
Because both liftings map into the same Hilbert space, their
difference defines a canonical linear operator.  The following theorem
is the conceptual turning point of the book: it shows that the entire
analytic development to come is already latent inside the WIS, waiting
only to be revealed by the logarithmic transformation.
\begin{theorem}[Canonical linear representation]\label{ch:thm:mismatch}
Let $A$ and $C$ be the lifting operators of
Proposition~\ref{ch:prop:canonical-lifting}.  The unique linear operator
$\Phi:V_M\oplus V_C\to V_I$ compatible with the logarithmic
factorisation $\log n_{mc}=v(c)-u(m)$ is
\[
\Phi(u,v)=Au-Cv,
\qquad\text{or in block form}\quad
\Phi=[A\;-C].
\]
\end{theorem}
\begin{proof}
The logarithmic factorisation is $\omega_{mc}=v(c)-u(m)$.
The canonical linear operator comparing lifted message and ciphertext
potentials is therefore $\Phi(u,v)=Au-Cv,$ whose edge values are $u(m)-v(c)$.
Consequently, $\omega = -\Phi(u,v),$ and equivalently $\omega=\Phi(-u,-v)$.
Thus the logarithmic factorisation determines the operator uniquely
up to the global sign convention fixed by the choice of mismatch.
\end{proof}

The mismatch operator is therefore not a modelling choice: it is the
canonical linear representation forced by the additive structure of
the logarithmic WIS representation.  Its image characterizes exact WIS
factorizability; compatibility with a prescribed prior is a separate
condition.
\begin{theorem}[Operator form of WIS factorisation]\label{ch:thm:operator-wis}
Let $\omega_{mc}=\log n_{mc}$.  The connected positive edge labels $n_{mc}$
admit a WIS factorisation if and only if
$\omega\in\operatorname{Im}(\Phi)$.  Equivalently, they satisfy the
cycle identities of Theorem~\ref{ch:thm:cycle-factorisation}.
\end{theorem}
\begin{proof}
A WIS factorisation exists iff
$\omega_{mc}=v_c-u_m$ for some $u\in V_M$, $v\in V_C$.  Since
$\Phi(u,v)=Au-Cv$, $\Phi(-u,-v) = -Au + Cv = v - u = \omega.$ Hence $\omega\in\operatorname{Im}(\Phi)$.  Conversely, every vector
in $\operatorname{Im}(\Phi)$ can, after changing the signs of its
vertex potentials, be written in the form $v_c-u_m$.  The equivalence
with the cycle identities is Theorem~\ref{ch:thm:cycle-factorisation}.
\end{proof}

Universal optimality relative to a fixed prior is stronger: in addition
to $\omega\in\operatorname{Im}(\Phi)$, the message potential in the
factorisation must satisfy $u=\log\pi$ up to an additive constant.
Thus a single linear operator $\Phi$ captures the entire
representability condition.  \textbf{From this point onward, every
analytic construction---the least-squares variational principle, the
normal operator, the posterior Laplacian, convexity, strict convexity
modulo gauge, stability of the reconstruction, the quotient geometry,
and the intrinsic distance from universal optimality---follows from
this canonical operator through standard Hilbert-space theory.}
Nothing is chosen or added; each step is forced by the previous one.
The graph-potential viewpoint developed in this section is explanatory
rather than an additional representation theorem; standard results on
graph potentials and incidence matrices provide the appropriate broader
context.  \paragraph{Running example: geo-indistinguishability.}
Geo-indistinguishability (Section~\ref{ch:sec:geo-indistinguishability-and-location-privacy}) provides a continuous analogue: the Laplacian noise mechanism produces a linear relation between the logarithm of the posterior density and the Euclidean distance to the reported location.  For two locations $P,Q$ at distance $r$, the log-incidence vector satisfies $\omega_{mc}=\log(\beta(c)/w(m))=\varepsilon r$ for the incident pair, so the least-squares solution of $\Phi(u,v)\approx-\omega$ reduces to a one-dimensional problem in the radial direction.

Its main contribution is to reveal the linear structure that
makes the operator-theoretic development of the next section possible.
\endinput
